\documentclass[12pt,a4paper]{ctexart}

\usepackage{amsmath,amssymb}
\usepackage{graphicx}
\usepackage{booktabs}
\usepackage{hyperref}
\usepackage{xcolor}
\usepackage{float}
\usepackage{enumitem}
\usepackage{url}

\graphicspath{{images/}}

\title{\heiti 复现 Grounding DINO：开放词汇目标检测与视觉定位}
\author{
    苏思齐 (12310645)，杨官宇涵 (12313614)，莫丰源 (12311805)\\
    南方科技大学 (SUSTech)，深圳\\[4pt]
    {\small\url{https://github.com/YangGuanyuhan/CV\_Project\_2026\_S}}
}
\date{\today}

\begin{document}

\maketitle

\begin{abstract}
开放词汇目标检测将传统检测扩展到任意文本描述的类别，使模型能够检测超出固定训练词汇表的物体。本报告复现了 Grounding DINO——一个最先进的开放词汇检测器，并在 COCO 2017 验证集上进行零样本评估。使用官方 Swin-Tiny 骨干网络权重，我们在完整的 val2017（5,000张图像）上取得了 AP=0.4055（AP50=0.5317）的结果。我们进一步进行了检测阈值敏感度和提示词格式设计的消融实验，发现点分隔的提示词格式对性能至关重要，且将 box\_threshold 从 0.35 降低到 0.25 可将子集 AP 从 0.4382 提升到 0.4637。成功和失败案例的定性分析表明，小目标检测和拥挤场景混淆是主要的局限性。
\end{abstract}

% ============================================================
\section{项目概述}
% ============================================================

\subsection{什么是开放词汇目标检测？}

传统目标检测系统（如 YOLO、Faster R-CNN）只能检测训练时定义好的固定类别（例如 COCO 的 80 类）。如果要检测新类别，就需要收集标注数据并重新训练模型——这是一个昂贵且耗时的工作。

\textbf{开放词汇目标检测}解决了这个问题：给定任意文本描述（如 "金色的狗 . 花盆 . 木椅子 ."），模型就能检测出对应的物体，不需要针对这些类别进行额外训练。

\subsection{Grounding DINO 简介}

Grounding DINO 是 2024 年 ECCV 发表的一个开放词汇检测模型，由 IDEA Research 开发。它结合了两个强大的技术：

\begin{itemize}
    \item \textbf{DINO}：一个基于 Transformer 的端到端目标检测器，用"目标查询"(object queries)替代传统的手工锚框
    \item \textbf{BERT}：一个预训练的语言模型，能将文本编码为上下文相关的向量表示
\end{itemize}

通过可变形注意力机制将视觉特征和文本特征融合，Grounding DINO 可以检测任意文本描述的物体。

\subsection{本项目做了什么}

我们使用官方发布的 \texttt{groundingdino-py} Python 包和 Swin-Tiny 预训练权重，完成以下工作：

\begin{enumerate}
    \item 在 COCO 2017 验证集（5,000张图像，80个类别）上进行完整的零样本评估
    \item 通过消融实验分析检测阈值和提示词格式对性能的影响
    \item 对检测成功和失败案例进行定性分析
    \item 构建模块化的 Python 推理、评估和可视化流水线
\end{enumerate}

\textbf{我们没有做}：从头训练、在 COCO 上微调、修改模型架构。

% ============================================================
\section{模型架构}
% ============================================================

\subsection{三大核心组件}

\begin{enumerate}
    \item \textbf{Swin Transformer 视觉骨干网络}
    \begin{itemize}
        \item 将输入图像处理为 4 个不同分辨率的层次化特征图
        \item 使用移位窗口注意力机制，计算复杂度与图像大小呈线性关系
        \item 配置：窗口大小=7，嵌入维度=96，各阶段深度=[2,2,6,2]
    \end{itemize}

    \item \textbf{BERT 语言编码器}
    \begin{itemize}
        \item 使用 bert-base-uncased 版本（12层 Transformer，隐藏维度768）
        \item 将文本提示词编码为上下文相关的词向量
        \item 最大文本长度：256 个 token
    \end{itemize}

    \item \textbf{可变形注意力融合模块}
    \begin{itemize}
        \item 6 层跨模态融合，每层 8 个注意力头
        \item 让 900 个目标查询同时关注视觉特征位置和文本 token 位置
        \item 每个查询输出：边界框、置信度分数、关联的文本短语
    \end{itemize}
\end{enumerate}

\subsection{推理流程}

\begin{enumerate}
    \item \textbf{加载模型}：加载 662MB 的预训练权重文件和架构配置
    \item \textbf{图像预处理}：BGR→RGB 转换，缩放到短边800px（长边不超过1333px），ImageNet 归一化
    \item \textbf{文本编码}：BERT 分词器将提示词转为 token，编码为 $(L, 768)$ 的文本特征矩阵
    \item \textbf{跨模态检测}：视觉特征与文本特征通过可变形注意力联合处理，输出检测结果
    \item \textbf{后处理}：归一化坐标 $(cx,cy,w,h)$ → 像素坐标 $(x_1,y_1,x_2,y_2)$ → COCO 格式 $(x,y,w,h)$
\end{enumerate}

% ============================================================
\section{实验设置}
% ============================================================

\subsection{数据集}

使用 \textbf{COCO 2017 验证集}：
\begin{itemize}
    \item 5,000 张 RGB 图像（室内、室外、街道等真实场景）
    \item 80 个常见物体类别（人、自行车、汽车、狗……）
    \item 36,781 个标注边界框（平均每张图约 7.4 个物体）
    \item 类别分布呈长尾形态：常见类如"人"有大量实例，稀有类如"吹风机"不到100个
\end{itemize}

\subsection{实验环境}

\begin{table}[H]
    \centering
    \caption{实验环境}
    \small
    \begin{tabular}{ll}
        \toprule
        \textbf{项目} & \textbf{配置} \\
        \midrule
        GPU & NVIDIA RTX 4060 Laptop, 8GB VRAM \\
        CPU & 13代 Intel Core i7, 16核 \\
        内存 & 32 GB \\
        Python & 3.10.20 \\
        PyTorch & 2.7.1 + CUDA 11.8 \\
        groundingdino-py & 0.4.0 \\
        \bottomrule
    \end{tabular}
\end{table}

\subsection{评估指标}

使用 COCO 标准指标（pycocotools COCOeval）：

\begin{itemize}
    \item \textbf{AP}：10 个 IoU 阈值 (0.50--0.95) 上的平均精度——主要指标
    \item \textbf{AP50}：IoU=0.50 时的 AP（宽松定位）
    \item \textbf{AP75}：IoU=0.75 时的 AP（严格定位）
    \item \textbf{APS / APM / APL}：按物体大小（小/中/大）分解的 AP
    \item \textbf{AR@k}：每张图取 k 个检测时的平均召回率
\end{itemize}

% ============================================================
\section{实验结果}
% ============================================================

\subsection{主要结果：完整 COCO val2017}

\begin{table}[H]
    \centering
    \caption{完整 COCO val2017 评估结果（5,000张图像）}
    \small
    \begin{tabular}{lclc}
        \toprule
        \textbf{精度指标} & \textbf{数值} & \textbf{召回指标} & \textbf{数值} \\
        \midrule
        AP & 0.4055 & AR@1 & 0.3273 \\
        AP50 & 0.5317 & AR@10 & 0.4756 \\
        AP75 & 0.4422 & AR@100 & 0.4875 \\
        APS (小) & 0.2587 & ARS (小) & 0.3156 \\
        APM (中) & 0.4328 & ARM (中) & 0.5154 \\
        APL (大) & 0.5510 & ARL (大) & 0.6472 \\
        \bottomrule
    \end{tabular}
\end{table}

\textbf{关键发现}：
\begin{itemize}
    \item 模型在 4,925/5,000 张图像中检测到物体（覆盖率 98.5\%），共产生 39,133 个检测
    \item 75 张图像（1.5\%）没有任何检测结果——通常是物体太小或遮挡严重的情况
    \item AR@100 (0.4875) > AP (0.4055)：模型能识别很多物体但定位不够精确，无法通过更严格的 IoU 阈值
    \item 大物体性能 (APL=0.551) 是小物体 (APS=0.259) 的约 \textbf{2.1 倍}
\end{itemize}

\subsection{消融实验 A：阈值敏感度}

测试了三种阈值设置（500张图像子集）：

\begin{table}[H]
    \centering
    \caption{阈值敏感度实验结果}
    \small
    \begin{tabular}{cccccc}
        \toprule
        \textbf{box\_th} & \textbf{text\_th} & \textbf{AP} & \textbf{AP50} & \textbf{APS} & \textbf{APL} \\
        \midrule
        0.25 & 0.20 & 0.4637 & 0.5884 & 0.3061 & 0.6338 \\
        0.35 & 0.25 & 0.4382 & 0.5532 & 0.2650 & 0.6269 \\
        0.45 & 0.30 & 0.3931 & 0.4827 & 0.2092 & 0.5841 \\
        \bottomrule
    \end{tabular}
\end{table}

\textbf{分析}：
\begin{itemize}
    \item 阈值越低，AP 越高，但检测数量也大幅增加
    \item 小物体对阈值最敏感：APS 从 0.209 提升到 0.306（+46\%）
    \item 大物体对阈值不敏感：APL 仅从 0.584 提升到 0.634（+8.5\%）
    \item 最佳阈值 box\_th=0.25 的 AP 比默认值高约 2.5 个百分点
\end{itemize}

\subsection{消融实验 B：提示词格式}

测试了三种提示词格式（500张图像子集）：

\begin{table}[H]
    \centering
    \caption{提示词格式对比}
    \small
    \begin{tabular}{lccc}
        \toprule
        \textbf{格式} & \textbf{示例} & \textbf{AP} & \textbf{AP50} \\
        \midrule
        点分隔 & \texttt{person . bicycle . car .} & 0.4382 & 0.5532 \\
        逗号分隔 & \texttt{person, bicycle, car,} & 0.1671 & 0.2038 \\
        句子风格 & \texttt{There are person, bicycle...} & 0.1675 & 0.2051 \\
        \bottomrule
    \end{tabular}
\end{table}

\textbf{核心发现}：点分隔格式比其他格式高出约 \textbf{62\% 的 AP}。原因在于 BERT 的分词器——点号周围有空格时被识别为独立的分隔 token，清晰划分了类别边界。逗号格式中逗号可能与前面的词合并（如 "car," 变成一个 token），句子格式中的虚词则引入了噪声。

\subsection{定性分析}

\subsubsection{成功案例}

模型能准确检测中大型物体，框与物体贴合紧密，类别标签正确。典型成功场景包括：户外街道（人、汽车、交通灯）、室内空间（椅子、桌子、杯子）、运动场景（人物、球拍、球）。

\subsubsection{失败模式}

\begin{enumerate}
    \item \textbf{小物体漏检（最常见）}：小于 32×32 像素的物体经常被漏掉。根本原因是 Swin-T 骨干网络的 32× 下采样导致小物体在特征图中几乎没有表示。
    \item \textbf{拥挤场景混淆（偶发）}：在物体密集的场景中（>15 个实例），模型可能产生重复框或类别混淆。
    \item \textbf{遮挡导致的定位误差（中等频率）}：部分遮挡降低检测置信度，并降低边界框精度。这解释了 AP50(0.53) 与 AP75(0.44) 之间的 9 个百分点的差距。
\end{enumerate}

% ============================================================
\section{讨论}
% ============================================================

\subsection{与官方结果的对比}

官方论文报告的 Swin-T 零样本 AP 约为 0.485，而我们得到的是 0.4055。这个约 8 个百分点的差距主要来自四个因素：

\begin{enumerate}
    \item 官方使用 GLIP 风格的评估协议（仅测试未见过的类别），我们使用标准 COCO 协议（全部 80 类）
    \item 官方可能不使用固定阈值，而我们使用了 box\_th=0.35 进行过滤
    \item 官方可能使用更丰富的提示词（类别描述、同义词），我们只用了标准类别名称
    \item \texttt{groundingdino-py} 包版本可能与原始代码有细微差异
\end{enumerate}

我们的结果（AP=0.4055）代表了使用固定阈值的标准 COCO 评估条件下更真实的零样本基线。

\subsection{主要局限性}

\begin{itemize}
    \item \textbf{小物体检测}：APS=0.259 vs APL=0.551，这是最关键的瓶颈
    \item \textbf{提示词格式敏感}：点分隔格式是强制要求，自然语言查询效果很差
    \item \textbf{计算成本}：完整评估需要约 68 分钟（0.81 秒/张），不适合实时应用
    \item \textbf{零检测图像}：1.5\% 的图像完全无检测输出
\end{itemize}

\subsection{未来方向}

\begin{enumerate}
    \item 在 COCO 训练集上微调（预期可将 AP 提升到约 0.506）
    \item 跨数据集评估（Objects365、OpenImages、LVIS）
    \item 与其他开放词汇检测器（GLIP、OWL-ViT、YOLO-World）的公平对比
    \item 自适应阈值策略以恢复零检测图像
    \item 模型蒸馏或 TensorRT 部署以降低推理延迟
\end{enumerate}

% ============================================================
\section{项目实现}
% ============================================================

\subsection{代码架构}

本项目实现为一个配置驱动的 Python 包（\texttt{cv-project-2026}），包含约 1,800 行源代码和 10 个 CLI 脚本：

\begin{table}[H]
    \centering
    \caption{项目模块组织}
    \small
    \begin{tabular}{llr}
        \toprule
        \textbf{模块} & \textbf{职责} & \textbf{代码行数} \\
        \midrule
        \texttt{src/models/grounding\_dino.py} & 模型包装器 & 182 \\
        \texttt{src/inference/predictor.py} & 推理接口 & 229 \\
        \texttt{src/inference/visualizer.py} & 可视化绘制 & 262 \\
        \texttt{src/inference/coco\_visualizer.py} & COCO 对比可视化 & 213 \\
        \texttt{src/engine/evaluator.py} & COCO 评估引擎 & 368 \\
        \texttt{src/datasets/coco\_categories.py} & 类别映射与提示词 & 269 \\
        \texttt{src/datasets/coco\_dataset.py} & 数据集工具 & 188 \\
        \texttt{src/utils/box\_ops.py} & 框操作 (IoU, NMS) & 115 \\
        \bottomrule
    \end{tabular}
\end{table}

\subsection{运行命令速查}

{\small
\begin{verbatim}
# 环境安装
conda create -n grounding_dino python=3.10 -y
conda activate grounding_dino
pip install torch torchvision --index-url https://download.pytorch.org/whl/cu118
pip install groundingdino-py
pip install -r requirements.txt

# 单张图像推理
python scripts/inference.py --image demo.jpg \
    --text "person . car . dog ." \
    --checkpoint checkpoints/groundingdino_swint_ogc.pth

# 完整 COCO 评估
python scripts/eval.py --config configs/grounding_dino.yaml \
    --checkpoint checkpoints/groundingdino_swint_ogc.pth \
    --coco_image_dir data/coco/val2017 \
    --coco_ann_file data/coco/annotations/instances_val2017.json

# 消融实验
python scripts/run_experiments.py \
    --checkpoint checkpoints/groundingdino_swint_ogc.pth \
    --coco_image_dir data/coco/val2017 \
    --coco_ann_file data/coco/annotations/instances_val2017.json \
    --subset_size 500
\end{verbatim}
}

% ============================================================
\section{结论}
% ============================================================

本项目成功复现了 Grounding DINO 的零样本开放词汇目标检测能力。核心成果总结：

\begin{enumerate}
    \item \textbf{基线性能}：在 COCO 2017 val2017（5,000张图像）上达到 AP=0.4055，AP50=0.5317
    \item \textbf{阈值影响}：AP 在 0.39--0.46 之间随阈值变化，最佳 box\_th=0.25
    \item \textbf{提示词格式至关重要}：非点分隔格式导致约 62\% 的 AP 下降
    \item \textbf{小物体是主要瓶颈}：APS=0.259 远低于 APL=0.551（2.1 倍差距）
    \item \textbf{定位优于分类}：AR@100=0.4875 > AP=0.4055，说明类别识别尚可但定位精度不足
\end{enumerate}

\subsection*{成员贡献}

\textbf{苏思齐 (12310645)}：文献调研、提示词消融实验设计与执行、失败案例定性分析。

\textbf{杨官宇涵 (12313614)}：整体项目设计与实现，包括全部 Python 模块和 CLI 脚本编写、完整 COCO 评估执行、消融实验执行、报告撰写。

\textbf{莫丰源 (12311805)}：阈值消融实验设计与执行、实验数据分析、可视化图表生成、报告实验部分撰写。

\end{document}
